% 声明为子文件，指定主文件
\documentclass[main.tex]{subfiles}

\begin{document}
\pagestyle{plain}
\setcounter{chapter}{2}

\chapter{Observed Regularities}
\label{chap:observed-regularities}


Beyond the classification, several quantitative regularities were noted across all trials.

\subsection{Internal Color Periodicity}

Along any directional line, the color pattern of heights repeats periodically without variation. For instance, a line with  $(2, 1)$ always follows the sequence (brown, white, brown) from left to right.

\subsection{Frequency of Occurrence}

Lines whose slope is close to horizontal or vertical appear less frequently across trials. The vast majority of observed lines are horizontal, vertical, or at a $45^\circ$ angle (slope~$\pm1$). Lines with slopes $1/2$, $1/3$, $2$, or~$3$ are comparatively rare.

\subsection{Vector Sum at Intersections}

When multiple directional lines meet at a single intersection point, the sum of their unit direction vectors equals zero. This suggests a form of local conservation governing the pattern geometry.


\end{document}